% MCQ Questions for Chapter 03
\section{বহুনির্বাচনি প্রশ্ন (MCQ)}

\subsection*{\textbf{\underline{সংখ্যা পদ্ধতি ও রূপান্তর}}}

\begin{multicols}{2}
\begin{enumerate}

\item দশমিক সংখ্যা পদ্ধতির ভিত্তি কত? \textbf{[ঢাকা-২০২৩]}
\begin{enumerate}[label=\alph*)]
\item ২
\item ৮
\item ১০ \checkmark
\item ১৬
\end{enumerate}

\item বাইনারি সংখ্যা পদ্ধতিতে কয়টি অঙ্ক ব্যবহৃত হয়? \textbf{[চট্টগ্রাম-২০২৩]}
\begin{enumerate}[label=\alph*)]
\item ২টি \checkmark
\item ৮টি
\item ১০টি
\item ১৬টি
\end{enumerate}

\item হেক্সাডেসিমেল সংখ্যা পদ্ধতিতে F এর মান কত? \textbf{[সকল বোর্ড-২০১৮]}
\begin{enumerate}[label=\alph*)]
\item ১০
\item ১৪
\item ১৫ \checkmark
\item ১৬
\end{enumerate}

\item (১০১০)$_২$ এর দশমিক মান কত? \textbf{[বরিশাল-২০২২]}
\begin{enumerate}[label=\alph*)]
\item ৮
\item ১০ \checkmark
\item ১২
\item ১৪
\end{enumerate}

\item (১৭)$_৮$ এর দশমিক মান কত? \textbf{[কুমিল্লা-২০২২]}
\begin{enumerate}[label=\alph*)]
\item ১৩
\item ১৪
\item ১৫ \checkmark
\item ১৬
\end{enumerate}

\item বাইনারি বিয়োগে কোন পদ্ধতিটি সবচেয়ে বেশি ব্যবহৃত হয়? \textbf{[যশোর-২০২২]}
\begin{enumerate}[label=\alph*)]
\item ২-এর পরিপূরক \checkmark
\item ১-এর পরিপূরক
\item সরাসরি বিয়োগ
\item কোনোটিই নয়
\end{enumerate}

\end{enumerate}
\end{multicols}

\subsection*{\textbf{\underline{লজিক গেট ও বুলিয়ান অ্যালজেব্রা}}}

\begin{multicols}{2}
\begin{enumerate}[resume]

\item বুলিয়ান অ্যালজেব্রায় ১ + ১ = কত? \textbf{[ঢাকা-২০২৩]}
\begin{enumerate}[label=\alph*)]
\item ০
\item ১ \checkmark
\item ১০
\item ১১
\end{enumerate}

\item সর্বজনীন (Universal) গেট কোনটি? \textbf{[যশোর-২০২৩]}
\begin{enumerate}[label=\alph*)]
\item NAND ও NOR \checkmark
\item AND ও OR
\item XOR ও XNOR
\item NOT ও AND
\end{enumerate}

\item XOR গেটের আউটপুট কখন ১ হয়? \textbf{[কুমিল্লা-২০২২]}
\begin{enumerate}[label=\alph*)]
\item ইনপুটগুলো ভিন্ন হলে \checkmark
\item ইনপুটগুলো একই হলে
\item সব ইনপুট ০ হলে
\item কোনোটিই নয়
\end{enumerate}

\item ডি মরগানের প্রথম সূত্র কোনটি? \textbf{[ঢাকা-২০২২]}
\begin{enumerate}[label=\alph*)]
\item (A+B)' = A'.B' \checkmark
\item (A.B)' = A'+B'
\item A+B = A.B
\item A.B = A+B
\end{enumerate}

\item ৩ ইনপুটবিশিষ্ট AND গেটের আউটপুট ১ হতে হলে ইনপুটগুলো হবে— \textbf{[চট্টগ্রাম-২০২৩]}
\begin{enumerate}[label=\alph*)]
\item ০, ০, ০
\item ১, ০, ১
\item ১, ১, ০
\item ১, ১, ১ \checkmark
\end{enumerate}

\end{enumerate}
\end{multicols}

\subsection*{\textbf{\underline{অ্যাডার ও ডিজিটাল সার্কিট}}}

\begin{multicols}{2}
\begin{enumerate}[resume]

\item হাফ অ্যাডার (Half Adder) এ আউটপুট কয়টি থাকে? \textbf{[ঢাকা-২০২৩]}
\begin{enumerate}[label=\alph*)]
\item ১টি
\item ২টি \checkmark
\item ৩টি
\item ৪টি
\end{enumerate}

\item এনকোডার (Encoder) এর কাজ কী? \textbf{[সিলেট-২০২৩]}
\begin{enumerate}[label=\alph*)]
\item মানুষের ভাষাকে কম্পিউটারের ভাষায় রূপান্তর \checkmark
\item কম্পিউটারের ভাষাকে মানুষের ভাষায় রূপান্তর
\item ডেটা ডিকোড করা
\item কোনোটিই নয়
\end{enumerate}

\item ডিকোডার (Decoder) এ n সংখ্যক ইনপুটের জন্য আউটপুট সংখ্যা কত? \textbf{[বরিশাল-২০২২]}
\begin{enumerate}[label=\alph*)]
\item $n$
\item $2n$
\item $2^n$ \checkmark
\item $n^2$
\end{enumerate}

\item রেজিস্টার (Register) তৈরিতে কী ব্যবহৃত হয়? \textbf{[দিনাজপুর-২০২৩]}
\begin{enumerate}[label=\alph*)]
\item লজিক গেট
\item ফ্লিপ-ফ্লপ \checkmark
\item এনকোডার
\item ডিকোডার
\end{enumerate}

\item কাউন্টার (Counter) এর প্রধান কাজ কী? \textbf{[কুমিল্লা-২০২২]}
\begin{enumerate}[label=\alph*)]
\item পালস গণনা করা \checkmark
\item ডেটা সংরক্ষণ
\item সিগন্যাল রূপান্তর
\item কোনোটিই নয়
\end{enumerate}

\end{enumerate}
\end{multicols}

\subsection*{\textbf{\underline{বহুপদী সমাপ্তিসূচক ও অভিন্ন তথ্যভিত্তিক প্রশ্ন}}}

\begin{multicols}{2}
\begin{enumerate}[resume]

\item \textbf{নিচের উদ্দীপকটি পড়ো এবং প্রশ্নের উত্তর দাও:}

\textit{রনি একটি ডিজিটাল সার্কিট তৈরি করল যার ৩টি ইনপুট এবং ৮টি আউটপুট লাইন রয়েছে।}

\item রনির তৈরি সার্কিটটি হলো—
\begin{enumerate}[label=\alph*)]
\item এনকোডার
\item ডিকোডার \checkmark
\item হাফ অ্যাডার
\item ফুল অ্যাডার
\end{enumerate}

\item বুলিয়ান অ্যালজেব্রার ক্ষেত্রে সঠিক তথ্য হলো—
\begin{enumerate}[label=\Roman*.]
\item এতে শুধু ০ এবং ১ ব্যবহৃত হয়
\item এটি লজিকের ওপর ভিত্তি করে কাজ করে
\item এতে কোনো জ্যামিতিক ক্যালকুলেশন নেই
\end{enumerate}

নিচের কোনটি সঠিক?
\begin{enumerate}[label=\alph*)]
\item i ও ii
\item i ও iii
\item ii ও iii
\item i, ii ও iii \checkmark
\end{enumerate}

\end{enumerate}
\end{multicols}
